% Source: https://artofproblemsolving.com/community/c3h2500010_my_300th_post_celebration
% Stable identity: aops-topic-2500010-q01
\textbf{(a)} Pair each number $k$ with $n+1-k$. Every pair has sum $n+1$, so each pair has average $\dfrac{n+1}{2}$; if $n$ is odd, the middle number $\dfrac{n+1}{2}$ is paired with itself and has the same average. A list made of groups that all have the same average has that average too, so the average of $1,\dots,n$ is $\dfrac{n+1}{2}$. The sum is the number of terms times the average:
\[1+2+\cdots+n=n\cdot\dfrac{n+1}{2}=\dfrac{n(n+1)}{2}.\]

\textbf{(b)} $1+2+\cdots+300=\dfrac{300\cdot301}{2}=\boxed{45150}$.
